\section{A reserved exponential window}\label{sec:coexistence}

This appendix uses the closed-code processor estimate of Appendix~\ref{app:code-estimates} (Lemma~\ref{x:exponential}), also used by Theorem~\ref{thm:exponential}, and shows that a decidable exponential-storage window can coexist in four dimensions with linear-storage windows whose storing-label set is undecidable.
Appendices~\ref{sec:fold-controls} and~\ref{sec:folded} supply the part governed by the halting computation.
The proof of Proposition~\ref{p:coexistence} uses the same internal-factor layer convention as
Section~\ref{s:model}.

\begin{proposition}[Coexistence of storage scales]\label{p:coexistence}
There is one computably specified translation-invariant, finite-range,
time-periodic circuit on \(\Z^4\), with a fixed site space of qubits
and elementary inventory
\(\{X,H,\mathrm{CNOT},\mathrm{Toffoli},R_0,D_Z,I\}\), under independent
single-qubit depolarization after every layer, with computable disjoint
rational windows \(I_0,I_1,\ldots\), such that:
\begin{enumerate}
\item On \(I_0\), one total computable rate-independent encoder and
intended-code decoders achieve entanglement fidelity at least \(15/16\)
for every integer \(0\le t\le2^{\lfloor c_{\exp}L\rfloor}\), for all
\(L\) above a computable threshold and one fixed rational \(c_{\exp}>0\).
\item On \(I_{n+1}\), if \(U(n)\uparrow\), Proposition~\ref{q:linear}
holds for label \(n+1\) of this circuit with its threshold
\(L_0(n+1)\); if \(U(n)\downarrow\), the all-encoding fidelity bound of
Theorem~\ref{f:forgetting} holds for the entire circuit, uniformly in
size and rate within the window.
\item There are computable constants \(A,a>0\), \(b\in\N\), an
integer \(C_0\ge2\), and \(L_{\rm EB}\ge C_0/2\) for this circuit,
all independent of \(j\). For every \(j\ge1\), every \(\eps\in I_j\)
and every integer \(L\ge L_{\rm EB}\), the entire channel has a
decomposition of the form \eqref{eq:u-eb} at round \(201L+C_0\).
The labels \(j\) for which
the optimal lifetime at fidelity \(7/8\) is unbounded in \(L\) at
every rate in \(I_j\) (Corollary~\ref{cor:window-undecidable}),
and the labels \(j\) for which the folded encoder and decoder of
Proposition~\ref{q:linear} have fidelity at least \(7/8\) at every
rational rate in \(I_j\), every integer \(L\ge L_0(j)\) and every
integer \(0\le t\le\lfloor cL\rfloor\), both equal
\(\{0\}\cup\{n+1:U(n)\uparrow\}\), a \(\Pi^0_1\)-complete set.
\end{enumerate}
Time counts rounds (Definition~\ref{def:round}). The first window is
reserved for a syntactically looping input.
Whole-site noise is not covered by this proposition.
\end{proposition}

The finite-torus memory estimate used by the auxiliary part is Lemma~\ref{x:exponential} in Appendix~\ref{app:code-estimates}. Local toric-code memory and sweep decoding have precedents in
Ahn \cite[Section~7.2 and Appendices~A\nobreakdash--B]{Ahn2004}, Kubica--Preskill
\cite{KubicaPreskill2019}, and the recent constructions of
Drechsler--Oreg--Aharonov and Ponarovsky \cite{DOA2026,Ponarovsky2026}.
That lemma is proved for the reflected-sweep circuit and the
fault support actually needed here.

\par\medskip\noindent\emph{Proof idea for Proposition~\ref{p:coexistence}.}
Reserve label zero for an explicit infinite loop. Add a separate toric-code memory whose cut lets the data through only when independent noise probes detect the top window. On the other windows the folded circuit supplies linear storage when appropriate, and the forgetting proof covers the whole register.

\begin{proof}[Proof of Proposition~\ref{p:coexistence}]
Rebuild the folded quantum construction of Appendices~\ref{sec:fold-controls}
and~\ref{sec:folded} with the fixed positive-input wrapper
\(V(1)=\text{an explicit infinite loop}\) and \(V(m+2)=U(m)\).
At rank \(a=j+n_0+1\), the scan calculation in
Section~\ref{sec:controller} calls \(V(j+1)\). Hence label zero
is the loop and label \(n+1\) simulates \(U(n)\) exactly.
Its decreasing nominal windows \(\mathcal J_j\) have rational gap
\[
 \delta_{\rm gap}=\inf \mathcal J_0-\sup \mathcal J_1>0,\qquad
 \vartheta=(\inf \mathcal J_0+\sup \mathcal J_1)/2.
\]

Add a separate closed four-dimensional toric-code quantum track
with the processor of Lemma~\ref{x:exponential}. There are no quantum
gates between its core and the folded core. All spatial quantum
work precedes the final cuts. Every retained auxiliary quantum
factor, including work and data-derived records, receives a local
acceptance-controlled cut. Clean rejection replaces it; clean
acceptance is identity. Disposable work resets before reuse.
The fixed influence catalogue and predecessor count \(\kappa\)
are chosen before detector depth or probe counts.

At each site and round use \(R+1\) disjoint banks of \(M\) reset
qubit probes, exposed to \(K_{\rm samp}\) sampling layers. Their
endpoint ORs \(E,D_1,\ldots,D_R\) are independent Bernoulli variables
with the refresh rate \eqref{eq:q-refresh}, independently across
site-rounds. Bank zero supplies the common semantic refresh.
The other banks define an exogenous detector
\[
 \mathsf D_{\rm test}=\one\{\textstyle\sum_{r=1}^R D_r\ge
                    \lfloor R\vartheta\rfloor+1\}.
\]
Choose \(R\) so that \(e^{-R\delta_{\rm gap}^2/2}\le\delta_D\).
Multiplying the centred Bernoulli moment bound from Hoeffding's lemma \cite[Lemma~1]{Hoeffding1963},
\(\mathbb E e^{\lambda(X-\mathbb EX)}\le e^{\lambda^2/8}\),
over the independent banks and optimizing gives
\[
 \sup_{\mathcal J_0}\Pr(\mathsf D_{\rm test}=0)\le\delta_D,\qquad
 \sup_{j\ge1,\,\mathcal J_j}\Pr(\mathsf D_{\rm test}=1)\le\delta_D.
\]
The integer threshold implements the strict comparison with
\(\vartheta\) exactly.

Let \(A\) contain all nonprobe sampling flags and all nonsampling
flags on every wire. It is independent of the complete endpoint
vector. The actual detector \(\mathsf D_{\rm act}\) equals \(\mathsf D_{\rm test}\) on
\(A=0\); on \(A=1\) its output may be arbitrary. Compute banks
serially with two OR bits and a binary test counter. Let \(S_{\rm np}\)
count all nonprobe semantic, mixed ancilla, counter and work wires;
it depends on \(R\), but not on \(M\). After choosing \(M\),
count the physical nonsampling depth \(J\), including both cuts, every
idle probe and every reset. Then
\[
 \nu\le[S_{\rm np}K_{\rm samp}+((R+1)M+S_{\rm np})J]\eps,
 \qquad
 \frac\nu\mu\le\frac{8S_{\rm np}}M+
       8\left(R+1+\frac{S_{\rm np}}M\right)\frac J{K_{\rm samp}}.
\]
Choose \(M\ge\max(S_{\rm np},32S_{\rm np}/\delta_{\rm ov})\), then
\(K_{\rm samp}\ge64(R+2)J/\delta_{\rm ov}\). The ratio is at most
\(3\delta_{\rm ov}/8\). At \(\eps=1/(2MK_{\rm samp})\),
\(\mu\ge1/16\), so all nominal windows below \(1/100\) have
preimages in the proven regime. The independence of the sampling banks uses independent single-qubit depolarization.

Choose the constants in this order. Fix both quantum catalogues,
\(p_{\rm aux}\) and \(\kappa\); choose the folded selector's
\(C_h,B\) and all its positive, negative and wrapped bounds. Put
\[
 \mu_{\rm cap}\le\min\{\text{folded cap},p_{\rm aux}/16,
                  1/[64(\kappa+1)],1/100\}.
\]
Choose \(n_0\) to place all nominal windows below that cap, fixing
the rational gap. Choose
\(\delta_D\le\min\{p_{\rm aux}/16,1/[64(\kappa+1)],1/64\}\)
and \(\delta_{\rm ov}\le1\) below the folded overlay margin. Choose \(R\),
then build the finite palette and all nonprobe hardware, counting
\(S_{\rm np}\). Next choose \(M\), form and count the actual
\(J\), and choose \(K_{\rm samp}\). Detector work is on site and changes depth without changing
the earlier influence catalogue. Finally enumerate rational inner
physical windows in \(\mu^{-1}(\mathcal J_j)\), exactly as in
Section~\ref{s:parameters}. Choose the coefficient \(c\) of
Proposition~\ref{q:linear} and its finite
chart thresholds after this entire circuit is fixed.

The main refresh swaps every semantic factor with its own mixed ancilla and disposes of old contents. Thus the folded part has the
same \(A,E,H=A\lor E\) source interface after recompilation, with
\(p\le2\mu\). It retains Proposition~\ref{q:linear}, including on
the reserved loop. On \(\mathcal J_0\), the auxiliary exogenous source is
\[
 Z=A\lor E\lor(1-\mathsf D_{\rm test}),\qquad
 \Pr(Z=1)\le\nu+\mu+\delta_D
                  \le2\mu_{\rm cap}+\delta_D< p_{\rm aux}.
\]
It is iid over physical site-rounds and independent of the quantum
input. On \(Z=0\) the auxiliary processor and accepting cut are
correct; all departures have the preselected bounded influence.
Lemma~\ref{x:exponential} proves part (i). Initialize the folded
part in any fixed allowed state and decode only the auxiliary part.

On every tail window use instead
\[
 B_{\rm open}=A\lor \mathsf D_{\rm test},\qquad
 \beta_{\rm open}\le\nu+\delta_D,
 \qquad\kappa\beta_{\rm open}<1/4.
\]
A zero of this iid field guarantees a clean rejecting cut of all
auxiliary quantum factors. A surviving path of \(s\) rounds crosses
\(s\) distinct site-round sources. Its probability, summed over the
at most \(\kappa^s\) paths, is at most \(4^{-s}\). Polynomial
unions give a small initial-site ancestor tail.

For a halting label combine this with the folded part's ancestor
tail, using one physical initial-site mask that retains a site whenever
either quantum part has an ancestor there. The two masks may be
correlated through \(A\). Under one joint history law,
the union mask has a small marginal and radius-\(O(s)\) determining
cones. Apply Lemma~\ref{f:domination} to this union, couple its two
disjoint copies, and sample two joint histories conditional
on those masks, as in Proposition~\ref{f:extension}. Route all
quantum factors at a retained initial site together. This gives an
extension of the entire circuit even for initial entanglement across
its two parts, hardware and reference. Iteration and the cloning bound of Lemma~\ref{f:bell} prove part (ii); finitely many smaller tori use simultaneous full
replacement with probability at least \((\inf I_j)^{NL^4}>0\).

On any tail window, the folded exceptional probability at
\(201L+C_0\) is polynomial times exponential. The auxiliary
\(L\)-round open-path exception, unioned over outputs, is at most
\(C'L^4\cdot4^{-L}\). Their union has the same exponential form.
On its complement all initial quantum dependence is gone and initial
autonomous metadata was measured. The entire good channel is therefore
measure-and-prepare. Its exogenous mixture proves part (iii)'s EB
bound and the sharp linear tail scale. This bound is not asserted
on the reserved window.

Finally choose the folded part's explicit linear-time encoder on
every label and ignore the auxiliary output in its decoder. It succeeds
on the loop and nonhalting tail; the whole-circuit halting bound gives
a strict finite counterexample on every halting tail label. The exact
rational evaluation and continuity argument of
Section~\ref{a:certification} applies to these seven-gate circuits.
Thus the certification and memory labels are precisely the stated
\(\Pi^0_1\)-complete set. Outside the selected windows the circuit's behaviour remains unclassified.
\end{proof}
